\documentclass{article}

\usepackage{amsmath,amssymb}
\usepackage{xurl}
\usepackage[hidelinks]{hyperref}
\setlength{\emergencystretch}{2em}

% Shared commands for notes in docs/paper-gaps/.

\DeclareUnicodeCharacter{2080}{\ifmmode{}_{0}\else\textsubscript{0}\fi}
\DeclareUnicodeCharacter{2081}{\ifmmode{}_{1}\else\textsubscript{1}\fi}
\DeclareUnicodeCharacter{2082}{\ifmmode{}_{2}\else\textsubscript{2}\fi}
\DeclareUnicodeCharacter{2099}{\ifmmode{}_{n}\else\textsubscript{n}\fi}

\newcommand{\Lean}{\textsc{Lean}}
\ExplSyntaxOn
\NewDocumentCommand{\leanid}{m}{
  \ifmmode
    \textsf{\detokenize{#1}}
  \else
    \group_begin:
    \urlstyle{sf}
    \tl_set:Nn \l_tmpa_tl {#1}
    \tl_replace_all:Nnn \l_tmpa_tl {\_} {_}
    \exp_args:NV \path \l_tmpa_tl
    \group_end:
  \fi
}
\ExplSyntaxOff
\providecommand{\tr}{\operatorname{tr}}
\newcommand{\tnleanIssueBase}{https://github.com/LionSR/TNLean/issues/}
\newcommand{\tnleanPrBase}{https://github.com/LionSR/TNLean/pull/}
\newcommand{\tnleanDocBase}{https://sirui-lu.com/TNLean/docs/}
\newcommand{\ghissue}[1]{\href{\tnleanIssueBase#1}{GitHub issue \##1}}
\newcommand{\ghpr}[1]{\href{\tnleanPrBase#1}{PR \##1}}
\newcommand{\doclink}[1]{\href{\tnleanDocBase#1.html}{\path{#1}}}

% Dirac notation and the identity operator, as in blueprint/src/macros/common.tex.
% \providecommand keeps note-local definitions authoritative.
\usepackage{dsfont}
\providecommand{\ket}[1]{|#1\rangle}
\providecommand{\bra}[1]{\langle #1|}
\providecommand{\braket}[2]{\langle #1 | #2 \rangle}
\providecommand{\ketbra}[2]{|#1\rangle\!\langle #2|}
\providecommand{\Id}{\mathds{1}}
\providecommand{\one}{\mathds{1}}
\providecommand{\C}{\mathbb{C}}
\providecommand{\MN}[1]{M_{#1}(\C)}
\providecommand{\MD}{M_D(\C)}

% External referencing for paper sources.  The build helper writes sanitized
% auxiliary files under build/paper-aux/ and excludes duplicated source labels.
\usepackage{xr}

% Import a generated paper auxiliary file with a caller-supplied label prefix.
% The candidate paths support builds from docs/paper-gaps/, the repository root,
% and build/paper-aux/ itself.
\newcommand{\importpaperaux}[2]{%
  \IfFileExists{../../build/paper-aux/#2.aux}{%
    \externaldocument[#1]{../../build/paper-aux/#2}%
  }{%
    \IfFileExists{build/paper-aux/#2.aux}{%
      \externaldocument[#1]{build/paper-aux/#2}%
    }{%
      \IfFileExists{#2.aux}{%
        \externaldocument[#1]{#2}%
      }{}%
    }%
  }%
}

% General external reference macros
\newcommand{\paperref}[2]{\ref{#1-#2}}
\newcommand{\papereqref}[2]{(\ref{#1-#2})}

% CPSV16 source (1606.00608 / MPDO-22-12-17-2.tex) external document and shortcuts
\importpaperaux{cpsv16-}{1606.00608/MPDO-22-12-17-2-tnlean}
\newcommand{\cpsvref}[1]{\paperref{cpsv16}{#1}}
\newcommand{\cpsveqref}[1]{\papereqref{cpsv16}{#1}}

% Verdict marker: \gapnote{<kind>}{<status>}. Kind is the policy
% classification (clarification, local-correction, scope-restriction,
% unfaithful, false-source, open-gap); status is open, wip (elimination
% underway), resolved, or historical. Machine-read by the site index;
% prints a verdict line.
\newcommand{\gapnote}[2]{\par\noindent\textbf{Verdict:} #1 (#2).\par}


\title{The Finite-Range Cyclic Knabe Coefficient}
\author{}
\date{2026-08-17}

\begin{document}
\maketitle
\gapnote{clarification}{resolved}

\section{Source statements}

For a nearest-neighbor chain, let $m$ be the number of local interaction
terms in each open window, and let $\gamma$ be the gap of the corresponding
open-window Hamiltonian. Knabe's finite-size criterion has threshold
\begin{align}
    \gamma
        &>\frac{1}{m}.
    \label{eq:knabe88_finite_range_coefficient_knabe_threshold}
\end{align}
This threshold originates in Ref.~\cite{Knabe1988Energy}. It is quoted for
matrix-product-state parent Hamiltonians at lines~1483--1489 of the local
source for Ref.~\cite{PerezGarcia2007Matrix} and at lines~2194--2197 of the
local source for Ref.~\cite{Cirac2021Matrix}. Those passages state only the
threshold.

The closed-form periodic coefficient in the nearest-neighbor case is
\begin{align}
    \delta
        &=\frac{m\gamma-1}{m-1}.
    \label{eq:knabe88_finite_range_coefficient_nearest_neighbor_coefficient}
\end{align}
The derivation below obtains this expression as the range-two specialization
of a finite-range coefficient.

\section{Finite-range derivation}

Let $N$ and $m$ be positive integers. Let $h_i$, for
$i\in\mathbb Z/N\mathbb Z$, be orthogonal projections, and define the
periodic Hamiltonian $H$, the length-$m$ window operators $A_s$, and the
separation-$q$ pair sums $Q_q$ by
\begin{align}
    H
        &=\sum_i h_i,
    \label{eq:knabe88_finite_range_coefficient_periodic_hamiltonian}\\
    A_s
        &=\sum_{a=0}^{m-1}h_{s+a},
    \label{eq:knabe88_finite_range_coefficient_window_hamiltonian}\\
    Q_q
        &=\sum_i(h_i h_{i+q}+h_i h_{i-q}).
    \label{eq:knabe88_finite_range_coefficient_pair_sum}
\end{align}
Assume
\begin{align}
    1
        &\leq R\leq m,
    \label{eq:knabe88_finite_range_coefficient_range_bounds}\\
    N
        &\geq2m,
    \label{eq:knabe88_finite_range_coefficient_cycle_size}
\end{align}
and assume that terms at cyclic separation at least $R$ commute. Cyclic
double counting gives
\begin{align}
    \sum_s A_s^2
        &=mH+\sum_{q=1}^{m-1}(m-q)Q_q.
    \label{eq:knabe88_finite_range_coefficient_window_square_sum}
\end{align}

For orthogonal projections, the operator inequality
\begin{align}
    h_i h_j+h_j h_i
        &\leq h_i+h_j
    \label{eq:knabe88_finite_range_coefficient_projection_pair_bound}
\end{align}
implies
\begin{align}
    Q_q
        &\leq2H.
    \label{eq:knabe88_finite_range_coefficient_pair_sum_upper_bound}
\end{align}
If $q\geq R$, the separated projections commute, so
\begin{align}
    Q_q
        &\geq0.
    \label{eq:knabe88_finite_range_coefficient_pair_sum_nonnegative}
\end{align}
Expansion of $H^2$ also gives
\begin{align}
    H+\sum_{q=1}^{m-1}Q_q
        &\leq H^2.
    \label{eq:knabe88_finite_range_coefficient_hamiltonian_square_lower_bound}
\end{align}

Define
\begin{align}
    c
        &=m-R+1.
    \label{eq:knabe88_finite_range_coefficient_baseline_coefficient}
\end{align}
Then
\begin{align}
    m-q
        &=c+(R-1-q).
    \label{eq:knabe88_finite_range_coefficient_weight_decomposition}
\end{align}
The correction term splits into short-range and separated contributions:
\begin{align}
    \sum_{q=1}^{m-1}(R-1-q)Q_q
        &=\sum_{q=1}^{R-2}(R-1-q)Q_q
          +\sum_{q=R}^{m-1}(R-1-q)Q_q.
    \label{eq:knabe88_finite_range_coefficient_correction_split}
\end{align}
The short-range contribution satisfies
\begin{align}
    \sum_{q=1}^{R-2}(R-1-q)Q_q
        &\leq
          2H\sum_{q=1}^{R-2}(R-1-q),
    \label{eq:knabe88_finite_range_coefficient_short_range_bound}
\end{align}
because its weights are positive and
(\ref{eq:knabe88_finite_range_coefficient_pair_sum_upper_bound}) holds
termwise. The separated contribution satisfies
\begin{align}
    \sum_{q=R}^{m-1}(R-1-q)Q_q
        &\leq0,
    \label{eq:knabe88_finite_range_coefficient_separated_bound}
\end{align}
because $R-1-q\leq0$ and $Q_q\geq0$. The nonnegativity applies throughout
the required range since
\begin{align}
    R
        &\leq q\leq m-1\leq N-R.
    \label{eq:knabe88_finite_range_coefficient_separation_range}
\end{align}
The triangular sum is
\begin{align}
    \sum_{q=1}^{R-2}(R-1-q)
        &=\frac{(R-1)(R-2)}{2}.
    \label{eq:knabe88_finite_range_coefficient_triangular_sum}
\end{align}

Substituting
(\ref{eq:knabe88_finite_range_coefficient_weight_decomposition}) into
(\ref{eq:knabe88_finite_range_coefficient_window_square_sum}), then applying
(\ref{eq:knabe88_finite_range_coefficient_hamiltonian_square_lower_bound}),
(\ref{eq:knabe88_finite_range_coefficient_short_range_bound}), and
(\ref{eq:knabe88_finite_range_coefficient_separated_bound}), yields
\begin{align}
    \sum_s A_s^2
        &\leq
          cH^2+(m-c+(R-1)(R-2))H.
    \label{eq:knabe88_finite_range_coefficient_preliminary_upper_bound}
\end{align}
The scalar coefficient simplifies as
\begin{align}
    m-c+(R-1)(R-2)
        &=(R-1)^2.
    \label{eq:knabe88_finite_range_coefficient_scalar_simplification}
\end{align}
Therefore
\begin{align}
    \sum_s A_s^2
        &\leq cH^2+(R-1)^2H.
    \label{eq:knabe88_finite_range_coefficient_window_square_upper_bound}
\end{align}

Assume that every window satisfies the local gap inequality
\begin{align}
    A_s^2
        &\geq\gamma A_s.
    \label{eq:knabe88_finite_range_coefficient_local_gap_inequality}
\end{align}
Summing over $s$ and using
\begin{align}
    \sum_s A_s
        &=mH
    \label{eq:knabe88_finite_range_coefficient_window_sum}
\end{align}
gives
\begin{align}
    H^2
        &\geq
          \frac{m\gamma-(R-1)^2}{m-R+1}H.
    \label{eq:knabe88_finite_range_coefficient_finite_range_bound}
\end{align}
The coefficient is positive under the condition
\begin{align}
    (R-1)^2
        &<m\gamma.
    \label{eq:knabe88_finite_range_coefficient_positive_coefficient_condition}
\end{align}
For $R=2$,
(\ref{eq:knabe88_finite_range_coefficient_finite_range_bound}) reduces to
the nearest-neighbor coefficient
(\ref{eq:knabe88_finite_range_coefficient_nearest_neighbor_coefficient}).

\section{Verdict}

The nearest-neighbor threshold
(\ref{eq:knabe88_finite_range_coefficient_knabe_threshold}) is attributed to
Ref.~\cite{Knabe1988Energy}. The arbitrary finite-range coefficient
(\ref{eq:knabe88_finite_range_coefficient_finite_range_bound}) is a TNLean
derivation, formalized by
\leanid{ProjectionGeometry.quadraticForm_sum_projections_of_cyclic_knabe}.
It is not attributed to Knabe, to
Theorem~2.1(i) of Nachtergaele's martingale argument
\cite{Nachtergaele1996Spectral}, or to the two MPS reviews
\cite{PerezGarcia2007Matrix,Cirac2021Matrix}.

The abstract result concerns a cyclic family of projections. Applying it
uniformly to sufficiently large MPS parent Hamiltonians still requires a
uniform gap for the corresponding nonwrapping open interaction.

\section{The Gosset--Mozgunov threshold}

For nearest-neighbor chains, Gosset and Mozgunov improve the threshold
(\ref{eq:knabe88_finite_range_coefficient_knabe_threshold}). Write $n=m+1$
for the number of sites of the open window and $\epsilon^\circ_N$ for the
gap of the periodic chain of $N$ sites. For $n>2$ and $N>2n$ they prove
\begin{align}
    \epsilon^\circ_N
        &\geq\frac56\,\frac{n^2+n}{n^2-4}
          \left(\gamma-\frac{6}{n(n+1)}\right),
    \label{eq:knabe88_finite_range_coefficient_gosset_mozgunov_bound}
\end{align}
so the local gap threshold becomes
\begin{align}
    \gamma
        &>\frac{6}{n(n+1)},
    \label{eq:knabe88_finite_range_coefficient_gosset_mozgunov_threshold}
\end{align}
which is strictly smaller than $1/(n-1)$ for all $n>3$ and has the optimal
order $n^{-2}$ \cite{Gosset2016Local}. Ref.~\cite{Cirac2021Matrix} cites
this improvement at line~2196 of its local source. The bound is Theorem~3,
in Section~2, of the arXiv version of Ref.~\cite{Gosset2016Local}; the
source is not tracked in this repository. The improvement is
unformalized: no declaration derives
Eq.~\ref{eq:knabe88_finite_range_coefficient_gosset_mozgunov_bound}, and
the bibliography entry \path{Gosset2016Local} in
\path{blueprint/src/references.bib} is cited by no blueprint chapter. Its
proof weights the terms of each window by a nonuniform coefficient
sequence, so it is not a specialization of the uniform-window derivation
above.

\bibliographystyle{alpha}
\bibliography{references}

\end{document}
